\label{sec:computational-method}

Вычислительные эксперименты предназначены для проверки внутренней логики M4,
поиска контрпримеров к слишком сильным интерпретациям и анализа
чувствительности. Они не оценивают производительность конкретной команды или
ИИ-инструмента: длительности, топологии и функции издержек в основном
синтетические. Все размеры выборок ниже относятся к полностью перечислимым
замороженным матрицам либо к воспроизводимому генератору, а не к случайной
выборке реальных проектов.

\subsection{Experimental Questions and Reporting Rules}

Каждая серия отвечает на заранее заданный вопрос и использует собственный
estimand. Таблица~\ref{tab:experiment-design} фиксирует эту связь до изложения
результатов и не позволяет смешивать доказанный оптимум, срок фиксированной
политики и наблюдение на конечной синтетической сетке.

\begin{table}[ht]
\centering
\scriptsize
\caption{Связь вычислительных серий с проверяемыми утверждениями}
\label{tab:experiment-design}
\begin{tabularx}{\textwidth}{@{}l>{\raggedright\arraybackslash}X>{\raggedright\arraybackslash}X>{\raggedright\arraybackslash}Xl@{}}
\toprule
Серия & Вопрос & Изменяемые факторы & Estimand и правило & Метод \\
\midrule
EXP-00 & Воспроизводятся ли варианты иллюстративного сценария? & Конфигурации 2--4 & $B_4$, $T^*$; сертификат $T^*=B_4$ & exact \\
EXP-01 & Всегда ли достижима граница $B_4$? & DAG, фазы, $P$ & $T^*/B_4-1$; доказанный положительный разрыв & exact \\
EXP-02 & Чем различаются доступная мощность и настроенная конфигурация? & $P$, $C(P)$, $\gamma(P)$, human-share & Профиль $T^*(P)$ & exact \\
EXP-03 & Идентифицируют ли агрегаты точный срок? & Фазовый профиль при равных агрегатах & Парная разность $T^*$ & exact \\
EXP-04--05 & Как гранулярность взаимодействует с $P$ и overhead? & Декомпозиция, $P$, тип и величина overhead & Оптимум по гранулярности и interaction contrast & exact \\
EXP-06 & Устойчивы ли масштабные и сравнительные выводы? & Общий масштаб и замороженные семейства ошибок & Инвариантность и сохранение ранжирования & exact/stress \\
EXP-07 & Переносятся ли направления эффектов на синтетические DAG? & $N$, топология, веса, размещение human-work & Paired contrasts для $T(\pi)$; отдельный exact-аудит & mixed \\
EXP-08 & Что меняет расположение фиксированного $H$? & Концентрация на пути, фазы, $P$, $\gamma/C$ & Парные изменения агрегатов, пути и $T(\pi)$; exact-аудит & mixed \\
\bottomrule
\end{tabularx}
\end{table}

\FloatBarrier

Во всех отчётах различаются четыре величины: структурная граница $B_4$,
доказанный оптимум $T^*$, найденное допустимое решение без доказательства
оптимальности и срок фиксированной политики $T(\pi)$. Обозначение $T^*$
используется только для запусков со статусом \texttt{OPTIMAL}. Допустимое
расписание длины $B_4$ доказывает достижимость границы, но совпадение нижней
оценки с лучшим найденным решением без закрытого solver gap недостаточно.

\subsection{Common Semantics and Experimental Quantities}

Во всех сериях действует одна семантика M4: начатая задача удерживает один
агентный поток до последней фазы, в том числе во время ожидания человека; все
человеческие фазы используют один непрерываемый ресурс мощности 1; преемник
становится доступен после завершения всех предшественников. Множитель $C(P)$
применяется к агентным фазам, $\gamma(P)$ --- к человеческим. Решатель
минимизирует только makespan; очередь и суммарная блокировка сообщаются как
производные характеристики, но не служат вторичной целью. Поэтому сообщаемое
$Q$ относится к одному возвращённому оптимуму, а не к минимуму $Q$ среди всех
makespan-оптимальных расписаний. В точной модели задача назначается при начале
первой фазы; фиксированная имитационная политика резервирует поток заранее при
выборе готовой задачи. Её очередь и срок относятся именно к этой политике.

Для серий с равной долей человеческой работы обозначим
\[
  r_h:=\frac{h_i}{k_i}.
\]
В канонических объектах EXP-01 и EXP-03 имеем $k_i=1$ и одинаковое $r_h$ у
всех задач; в иллюстративном объекте EXP-02 сохраняется исходное $k_i$, а
$h_i=r_hk_i$. Профиль \texttt{io} делит $H_i$ пополам вокруг единственной
agent-фазы: $H_i/2,A_i,H_i/2$. Профиль \texttt{front\_loaded} использует
$0.9H_i,A_i,0.1H_i$. В профиле с равным чередованием четыре human-фазы имеют
длины $H_i/4$, а три agent-фазы --- $A_i/3$. Историческая метка
\texttt{alternating\_sync} обозначает именно это внутризадачное равенство и
\emph{не} навязывает совпадение запросов разных задач по wall-clock. В
\texttt{alternating\_staggered} human-фазы те же, а agent-доли
$\{1/6,2/6,3/6\}$ циклически переставляются между задачами. В EXP-08
нейтральное имя \texttt{equal\_alternating} используется для равного
чередования. Только в EXP-08 номинальные половины, четверти и трети
реализуются на allocation-сетке 0.001: целое число тиков делится с остатком,
который распределяется по seeded-перестановке. Сумма фаз сохраняется точно, а
части одного типа могут различаться не более чем на один тик.

Экспериментальная гранулярность $m$ --- число параллельных подзадач, которыми
заменяется выбранная задача. В каноническом операторе EXP-04--EXP-05 её
исходная входная human-фаза и полезная agent-работа делятся поровну между $m$
подзадачами; все они наследуют предшественников исходной задачи. После них
создаётся join, содержащий исходную финальную приёмку, сохраняющий исходный
quality gate и становящийся предшественником всех прежних successors. При
$m=1$ это семантически тот же объём полезной работы, хотя DAG уже содержит
явный join. Поэтому contrasts EXP-04--EXP-05 относятся к преобразованному
оператору «части плюс join»; C3 в EXP-07 использует исходную задачу при $m=1$, и
численные величины этих двух серий напрямую не сопоставляются.

Пусть $E_0$ --- исходная полная длительность выбранной задачи, $r_o$ --- ставка
издержек одной дополнительной границы, а $\lambda_h$ --- человеческая доля
этих издержек. Тогда
\[
  O(m,r_o)=(m-1)r_oE_0,\qquad
  O_h=\lambda_hO,\qquad O_a=(1-\lambda_h)O.
\]
Половина $O_h$ и $O_a$ поровну добавляется ко входам подзадач, вторая половина
--- к join. При $r_o=0$ полезные $A$ и $H$ сохраняются точно. Для двух уровней
параллелизма complementarity декомпозиции измеряется contrast
\begin{equation}
  I(P,m)=\bigl[T^*(P,1)-T^*(P,m)\bigr]
  -\bigl[T^*(1,1)-T^*(1,m)\bigr].
  \label{eq:interaction-contrast}
\end{equation}
Положительное $I$ означает, что выигрыш от дробления при $P$ больше, чем при
$P=1$. Contrast вычисляется только когда все четыре objectives доказаны как
\texttt{OPTIMAL}. Для заранее фиксированной политики $\pi$ отдельно обозначим
\[
  I_{\pi}(P,m)=\bigl[T_{\pi}(P,1)-T_{\pi}(P,m)\bigr]
  -\bigl[T_{\pi}(1,1)-T_{\pi}(1,m)\bigr],
\]
где $T_{\pi}$ --- срок этой политики, а не доказанный оптимум. В замороженных
матрицах EXP-04--EXP-05 используется
$\mathrm{tolerance}=10^{-4}$: contrast считается положительным при
$I>\mathrm{tolerance}$, отрицательным при $I<-\mathrm{tolerance}$ и нулевым в
остальных случаях. На конечной сетке EXP-04 кривая называется U-образной, если
все её точки доказаны оптимальными, минимум достигается хотя бы при одном
внутреннем $m\notin\{1,8\}$ и оба края $T^*(1)$ и $T^*(8)$ превосходят минимум
больше $\mathrm{tolerance}$. Это операциональное определение не утверждает U-образность
вне исследованных целых $m$.

Для синтетических DAG используются две концентрации. Величина
$\rho_L=L/W_4$ задаёт долю структурной работы на longest path при
$C=\gamma=1$, а $\rho_H^{cp}=H_{cp}/H$ --- долю всего human-work,
расположенную на выбранном структурном пути. Они не взаимозаменяемы: первая
характеризует последовательность DAG, вторая --- размещение фиксированного или
приближённо фиксированного человеческого ресурса.

\subsection{Frozen Exact Matrices EXP-00--EXP-05}

EXP-00 решает три полностью заданных варианта иллюстративного сценария
(варианты 2--4). EXP-01 использует полный факторный дизайн
\[
  5\ \text{топологий}\times N{=}4\times
  P\in\{2,4\}\times r_h\in\{0.10,0.25,0.50\}\times
  3\ \text{профиля}=90
\]
точных точек. Топологии --- chain, independent, fork--join, diamond и
two-layer; веса задач циклически равны $6,12,18,24$. Исходная замороженная
матрица содержала также $N=8,12$ (270 точек), но заранее заданное правило
требовало глобально уменьшить размер, если хотя бы одна точка не получает
\texttt{OPTIMAL} за общий 60-секундный лимит. Первая такая точка при $N=8$
получила \texttt{FEASIBLE}, после чего без индивидуальной настройки лимитов
была исполнена полная $N=4$-матрица; исходная матрица и запись решения
сохранены.

EXP-02 содержит два объекта --- DAG иллюстративного варианта 4 и канонический
fork--join при $N=4$ --- и сетку
$P\in\{1,2,3,4,6,8,12\}$. Для каждого объекта исполняются пять уникальных
режимов: три значения $r_h\in\{0.10,0.25,0.50\}$ при $C=\gamma=1$, а также
при $r_h=0.25$ отдельно
\[
 C(P)=1+0.05(P-1)+0.005P(P-1)
 \quad\text{и}\quad
 \gamma(P)=1+0.10(P-1).
\]
Общий режим $r_h=0.25,C=\gamma=1$ не дублируется, поэтому всего решается
$2\times5\times7=70$ точек.

EXP-03 скрещивает два DAG (fork--join и two-layer),
$P\in\{2,4\}$, $r_h\in\{0.25,0.50\}$ и четыре профиля
(\texttt{io}, \texttt{front\_loaded}, равное и переставленное чередование):
всего 32 точные задачи. Для каждого из восьми наборов
``DAG--$P$--$r_h$'' три профиля сравниваются с \texttt{io}, что даёт 24
парных contrasts при одинаковых $A,H,W_4,L_4,B_4,C,\gamma$.

EXP-04A повторно использует доказанную иллюстративную пару вариантов 3--4.
EXP-04B применяет описанный оператор к sink-задаче $E_0=24$ канонического
two-layer DAG при $N=4$, $P=4$, $r_h=0.25$ и профиле \texttt{io}. Сетка
\[
 m\in\{1,2,3,4,6,8\},\quad
 r_o\in\{0,0.025,0.05,0.10,0.20\},\quad
 \lambda_h\in\{0,0.5,1\}
\]
содержит 90 формальных комбинаций и 66 семантически различных сценариев:
при $m=1$ overhead нулевой для всех $r_o,\lambda_h$, а при $r_o=0$ значения
$\lambda_h$ совпадают. После восстановления общих точек анализируется 13
содержательных кривых.

EXP-05 использует тот же объект и оператор, но скрещивает
$P,m\in\{1,2,3,4,6,8\}$, $r_o\in\{0,0.05,0.10\}$ и
$\lambda_h\in\{0,0.5\}$. Из 216 формальных позиций получено 156 уникальных
задач; их решения восстанавливают 180 аналитических позиций после объединения
тождественных нулевых срезов. Первый общий 120-секундный проход дал 155
\texttt{OPTIMAL} и один \texttt{FEASIBLE}; предусмотренное матрицей единственное
увеличение лимита до 300 секунд было затем одинаково применено ко всем 156
точкам, а не только к трудной.

\subsection{Robustness, Synthetic DAGs, and Confirmatory Freeze}

EXP-06 разделён на три части. EXP-06A умножает четыре заранее выбранных объекта
(иллюстративные варианты 3--4 и две канонические точки) на
$\kappa\in\{0.25,0.5,1,2,4\}$, одновременно масштабируя $X$, фазы, time unit и
tolerance; это 20 точных точек с неизменной целочисленной моделью. EXP-06B
сравнивает варианты 3--4 при пяти уровнях
$u\in\{0.05,0.10,0.20,0.30,0.50\}$ и 1000 seeds на уровень: 5000 пар, или
10000 точных задач. Agent- и human-множители выбираются независимо и равномерно
в log-space на $[\log(1-u),\log(1+u)]$; общие задачи получают одинаковые
множители в двух вариантах, а декомпозированные тестовые задачи наследуют
множитель исходной задачи с независимым ограниченным log-residual масштаба
0.25. EXP-06C проверяет два заранее заданных неблагоприятных направления на
$u=0,0.005,\ldots,0.95$: исходная тестовая задача варианта 3 получает
$1-u$, а соответствующие декомпозированные задачи варианта 4 --- $1+u$ на
выбранном agent- либо human-ресурсе. Получаются 382 пары, или 764 точные
задачи. Эти семейства ошибок синтетические и не являются доверительными
интервалами эмпирической калибровки.

EXP-07 --- стратифицированный пилот из 72 ячеек:
\[
\begin{aligned}
N &\in \{8,16,32,64\}, &
\text{topology} &\in \{\text{wide-layered,mixed,chain-like}\},\\
\text{weights} &\in \{\text{homogeneous,lognormal}\}, &
\rho_H^{cp} &\in \{\text{low,medium,high}\}.
\end{aligned}
\]
В каждой ячейке принимаются 20 последовательных observation seeds, всего 1440
сценариев: 480 общих структурных основ «граф плюс веса» повторяются для трёх
уровней размещения human-work. Отклонённые попытки сохраняются, но наблюдениями
не считаются. Три
семейства топологии по построению соответствуют low, medium и high корзинам
$\rho_L$, поэтому эффект topology family не идентифицируется отдельно от
$\rho_L$. На полной выборке четыре замороженных направления C1--C4 оценивают
фиксированную политику или ветви $B_4$: различие равного и переставленного
чередования; сдвиг policy-optimal $P$ на сетке $\{1,2,4,8\}$ при росте
agent- или human-cost; policy-contrast $I_\pi(4,3)$ для
$m=3$, $P=1/4$, $r_o=0/0.10$, $\lambda_h=0.5$; и частоту активности human-
ветви при разных $\rho_H^{cp}$. При равенстве минимальных сроков policy-optimal
параллелизм определяется как
$P_{\mathrm{opt}}:=\min\operatorname*{arg\,min}_{P}T_\pi(P)$.

Знаки и правила решения были зафиксированы до расчёта. Для C1 ожидается
неотрицательная медиана разности policy-tightness равного и переставленного
чередования; решение требует неотрицательных point estimate и верхней границы
95\%-го bootstrap-интервала. Для C2 ожидается неположительный медианный сдвиг
policy-optimal $P$; кроме этого доля положительных сдвигов не должна превышать
0.05. Для C3 ожидаются неотрицательные медианы $I_\pi(4,3)$ и ослабления
выигрыша при переходе от $r_o=0$ к $r_o=0.10$. Для C4 ожидаются положительные
парная средняя разности индикаторов human-active branch и разность её частот
между high и low концентрациями. Индивидуальные исключения сохраняются и
сообщаются, поэтому подтверждение медианного направления не означает
универсальный знак на каждом DAG. Для интервалов используются 2000 кластерных
bootstrap-перевыборок с seed 707071; единицей перевыборки служит общая
структурная основа, чтобы три размещения human-work не считались независимыми.

После 20 принятых seeds в каждой ячейке замороженная матрица предписывала
pilot-extension rule для двух ключевых ячеек. В первой $N=16$, topology
family --- mixed, веса --- lognormal и $\rho_H^{cp}$ --- medium; во второй
$N=64$ при тех же topology family и типе весов, но $\rho_H^{cp}$ --- high.
По протоколу выборка должна была увеличиться до 50 seeds, если
bootstrap-интервал первичной медианы пересекает ноль либо Monte Carlo standard
error первичной доли превышает 0.05. Реализованный runner проверил только
пересечение нуля глобальными bootstrap-интервалами медиан; ветви key-cell и
MCSE не были исполнены. Глобальный триггер не сработал, extension не запускался,
поэтому EXP-07 следует интерпретировать как 20-seed pilot с частично
реализованным правилом расширения, а не как полностью завершённую
confirmatory-процедуру. Заранее выбранный exact-аудит содержит 12 DAG с
$N\in\{8,16\}$, homogeneous-весами, средней концентрацией и seeds 0--1 для
трёх topology families. После единой ресурсной поправки, принятой до просмотра
claim outcomes, exact-scope ограничен baseline-gap и двумя профилями C1: 36
сценариев с лимитом 10 секунд. C2 и C3 в EXP-07 являются policy-only, C4 ---
lower-bound-only. \texttt{FEASIBLE} incumbents исключаются из доказанных пар.

EXP-08 исправляет смешение общего $H$ и его расположения с помощью независимых
development и confirmatory потоков. Общая сетка содержит 24 ячейки
$N\times\text{topology}\times\text{weight type}$. Development namespace даёт
по 10 основ на ячейку (240), confirmatory namespace --- по 20 (480); raw seed
равен первым 64 битам SHA-256 от namespace, ячейки, номера наблюдения и номера
попытки. Confirmatory-графы генерируются только при наличии freeze-файла, а
пересечения raw seeds с development и EXP-07 автоматически проверяются.

Для каждой основы сначала при $C=\gamma=1$ выбирается reference path; равные
longest paths разрешаются seeded-рангом, не зависящим от task id, а их
множественность сохраняется отдельным признаком. Общая доля human-work строго
фиксируется как $H/W_4=0.12$, после чего на reference path размещается ровно
0.30 или 0.70 от $H$. Распределение внутри и вне пути пропорционально
структурным весам, ограничено task-level долями $[0.025,0.85]$ и выполняется
bounded largest-remainder allocator на сетке 0.001 с seeded tie-break. Для
schedule-effect используются профили \texttt{io} и
\texttt{equal\_alternating}, $P\in\{1,4,8\}$ и $C=\gamma=1$; парная разность
``high минус low'' нормируется на общий $B_4$. Engineering equivalence margin
равен 0.01, sensitivity margins --- 0.005 и 0.02; 2000 bootstrap-перевыборок
выполняются по graph base с seed 808071. Отдельный differentiated-slowdown
срез использует $C=1$, $\gamma=1.3$, $P=4$.

Scope exact-аудита EXP-08 выбран только по statuses и wall time development-
прогонов, без сохранения их objectives. После экранов с общими лимитами 5 и 30
секунд до confirmatory-генерации были заморожены $N=8$, homogeneous-веса,
семейства mixed/chain-like, observation indices 0--1, $P\in\{4,8\}$ и обе
allocations. Это 16 confirmatory-сценариев с лимитом 30 секунд; wide-layered и
$N=16$ исключены единообразно по timing rule, а не по значению эффекта.

\subsection{Synthetic DAG Generator}

EXP-07 и EXP-08 используют один тип слоистого генератора. Число слоёв равно
округлённому $Nf$ с ограничением $[2,N]$. При двух слоях source образует первый
слой, а все остальные вершины --- второй; при трёх и более слоях первый и
последний слои содержат по одной вершине, а остальные вершины распределяются
между промежуточными слоями.
Каждая вершина получает хотя бы одного предшественника из соседнего слоя, и
каждая вершина предыдущего слоя получает хотя бы одного successor. Параметры
$(f,p_{adj},p_{skip})$ равны $(0.20,0.30,0)$ для wide-layered,
$(0.45,0.25,0.04)$ для mixed и $(0.75,0.15,0)$ для chain-like. Принимаются
только ацикличные, слабо связанные графы, достижимые из source, с
$\rho_L\in[0.10,0.30)$, $[0.30,0.60)$ или $[0.60,0.90]$ соответственно.

Homogeneous-веса равны 12. Для lognormal-весов сначала генерируется
$Y\sim\operatorname{LogNormal}(0,0.55)$, затем множитель
$\min\{6,\max\{1,\operatorname{round}(2Y)\}\}$ умножается на 12. В EXP-07
две task-level human-share величины --- на структурном critical path и вне его
--- выбираются с сетки $0.025,0.05,\ldots,0.80$: сначала минимизируется
отклонение $\rho_H^{cp}$ от целей 0.18, 0.45 или 0.75 в соответствующей
корзине, затем отклонение общей human-share от 0.25. Это приближённое, а не
строгое выравнивание общего $H$; строго фиксированный $H$ вводится только в
EXP-08.

\subsection{Exact Solver, Time Grid, and Fixed Policy}

Точные задачи сформулированы как интервальная модель CP-SAT из OR-Tools
9.15.6755 \cite{perron2023cpsatlp}. Для фаз задаются целочисленные времена
начала и окончания; порядок фаз и дуги DAG задаются ограничениями
предшествования, человеческий ресурс --- глобальным \texttt{NoOverlap}, а
закрепление задачи и локальная блокировка --- альтернативными интервалами,
охватывающими весь span задачи на выбранном агентном потоке. Для каждого потока
также задаётся \texttt{NoOverlap}; objective --- максимум окончаний последних
фаз.

Каждый сценарий задаёт положительный \texttt{time\_unit}. Эффективная
длительность каждой фазы обязана делиться на него без остатка; иначе scenario
validator завершает работу ошибкой. Поэтому CP-SAT получает точную
целочисленную модель, а не округляет непрерывные длительности внутри решателя.
Обратное преобразование умножает integer ticks на \texttt{time\_unit}. В
EXP-06B/C возмущённые фазы заранее квантуются методом
\texttt{ROUND\_HALF\_UP} минимум до одного тика, после чего последняя фаза
каждого ресурса корректируется так, чтобы итог задачи оставался точно
представимым. EXP-08 отдельно распределяет $H$ на allocation-сетке 0.001, а
решает расписание на сетке 0.0001.

\begin{table}[ht]
\centering
\small
\caption{Целочисленные сетки и лимиты точного решателя}
\label{tab:solver-limits}
\begin{tabularx}{\textwidth}{@{}l>{\raggedright\arraybackslash}X>{\raggedright\arraybackslash}X@{}}
\toprule
Серия & Time unit & Общий лимит на сценарий \\
\midrule
EXP-00 & 0.05 & 300 s \\
EXP-01 & 0.025 & 60 s \\
EXP-02 & 0.00005 & 60 s \\
EXP-03 & 0.025 & 60 s \\
EXP-04 & 0.00625 & 120 s \\
EXP-05 & 0.00625 & 120 s, затем единое замороженное увеличение до 300 s \\
EXP-06A & исходная сетка объекта, умноженная на $\kappa$ & 60 s \\
EXP-06B/C & 0.05 & 10 s \\
EXP-07 exact-аудит & 0.001 & 10 s \\
EXP-08 exact-аудит & 0.0001 & 30 s \\
\bottomrule
\end{tabularx}
\end{table}

Все точные запуски используют один search worker и solver seed 0. Статусы,
objective, best bound, относительный gap и wall time сохраняются. Ограничение
времени не является доказательством оптимальности: строки \texttt{FEASIBLE}
сохраняют incumbent, но не получают обозначение $T^*$.

Для больших выборок используется заранее зафиксированная детерминированная
политика \texttt{bottom\_level\_fcfs\_v1}. Пусть $w_i^{(4)}$ --- эффективный
вес задачи, а
\[
  b_i=w_i^{(4)}+\max_{j\in\operatorname{succ}(i)}b_j
\]
--- её bottom level, где максимум пустого множества равен нулю. В каждый момент
события готовые неназначенные задачи сортируются по убыванию $b_i$, затем по
убыванию $w_i^{(4)}$ и лексикографически по task id; свободные потоки берутся
по возрастанию номера. Завершившаяся задача освобождает поток только после
последней фазы. Запросы к человеку обслуживаются по времени поступления
(FCFS); при одинаковом времени используются убывающий $b_i$, task id и индекс
фазы. При одном timestamp сначала продолжаются следующие фазы уже назначенных
задач, затем назначаются вновь готовые задачи, после чего запускается первый
human-запрос. Политика непрерываема и не использует будущую информацию сверх
фиксированного DAG и длительностей. Её результаты относятся к $T(\pi)$, а не к
$T^*$; небольшие точные подвыборки служат аудитом и не подменяют основной
estimand.

\subsection{Reproducibility Environment and Integrity}

Матрицы EXP-01--EXP-08, сценарии EXP-00, правила принятия решений и seeds
хранятся вместе с кодом. Manifests записывают SHA-256 входов, реализации и
выходных файлов, версии зависимостей, solver settings и статусы. Пакет требует
Python $\ge3.11,<3.14$; команды воспроизведения запрашивают Python 3.13, а
\texttt{uv.lock} и \texttt{pyproject.toml} фиксируют OR-Tools 9.15.6755,
PyYAML 6.0.3 и Matplotlib 3.11.1. Сохранённое локальное окружение, на котором
проверялась эта редакция, сообщает Python 3.13.14 и macOS 26.5 arm64.

При этом исходные manifests не сохраняют patch-версию Python, ОС, процессор или
объём памяти. Поэтому указанные Python patch и платформа характеризуют
доступную локальную проверку, но не являются доказанными метаданными каждого
первичного прогона; wall time нельзя воспроизводить как платформенно
инвариантную величину. Численные выводы основаны на статусах и objectives, а не
на сравнении wall time. Команды, порядок EXP-08 и полный список артефактов
приведены в Приложении~\ref{app:experiments}.
